\documentclass[12pt]{article}
\usepackage{amsmath,amssymb,amsthm}
\usepackage[margin=1in]{geometry}
\usepackage{algorithm}
\usepackage{algpseudocode}

\newtheorem{theorem}{Theorem}
\newtheorem{lemma}[theorem]{Lemma}

\title{Problem 205: Dice Game}
\author{Project Euler Solution}
\date{}

\begin{document}
\maketitle

\section{Problem Statement}
Peter has nine 4-sided dice (faces 1--4). Colin has six 6-sided dice (faces 1--6).
They each roll all their dice and compare totals. Find $\Pr(\text{Peter beats Colin})$,
rounded to 7 decimal places.

\section{Mathematical Foundation}

\begin{theorem}[Convolution of Discrete Uniform Distributions]
Let $X_1, \ldots, X_n$ be independent, each uniform on $\{1, \ldots, d\}$. The number of
outcomes giving sum $S_n = s$ is the coefficient of $x^s$ in $(x + x^2 + \cdots + x^d)^n$.
\end{theorem}

\begin{proof}
For independent discrete variables, $P(S_n = s) = \sum_{k=1}^{d} P(S_{n-1} = s-k) \cdot P(X_n = k)$,
which is the convolution of the individual PMFs. The generating function representation
follows from the product rule for GFs of independent variables.
\end{proof}

\begin{theorem}[Winning Probability]
Let $P$ denote Peter's total and $C$ Colin's total. Then:
\[
\Pr(P > C) = \frac{1}{4^9 \cdot 6^6}
\sum_{p=9}^{36} f_P(p) \cdot F_C(p-1),
\]
where $f_P(p)$ is the number of outcomes giving Peter total $p$, and
$F_C(t) = \sum_{c=6}^{t} f_C(c)$ is Colin's cumulative frequency.
\end{theorem}

\begin{proof}
Since $P$ and $C$ are independent:
\[
\Pr(P > C) = \sum_{p=9}^{36} \Pr(P=p) \cdot \Pr(C < p)
= \sum_{p=9}^{36} \frac{f_P(p)}{4^9} \cdot \frac{F_C(p-1)}{6^6}.
\]
Factoring out the denominators yields the formula.
\end{proof}

\begin{lemma}[Range Constraints]
Peter's sum ranges over $[9, 36]$ with $4^9 = 262144$ total outcomes.
Colin's sum ranges over $[6, 36]$ with $6^6 = 46656$ total outcomes.
\end{lemma}

\begin{proof}
The minimum sum of $n$ dice each showing at least~1 is $n$; the maximum with $d$-sided
dice is $nd$. For Peter: $[9, 36]$. For Colin: $[6, 36]$.
\end{proof}

\section{Algorithm}

\begin{algorithmic}[1]
\State Build frequency table $f_P[0\ldots 36]$ via 9-fold convolution of the uniform distribution on $\{1,2,3,4\}$
\State Build frequency table $f_C[0\ldots 36]$ via 6-fold convolution of the uniform distribution on $\{1,\ldots,6\}$
\State Compute cumulative $F_C[t] = \sum_{c=6}^{t} f_C[c]$ for $t = 6, \ldots, 36$
\State $\mathrm{win} \gets \sum_{p=9}^{36} f_P[p] \cdot F_C[p-1]$
\State \Return $\mathrm{win} \;/\; (4^9 \cdot 6^6)$, rounded to 7 decimal places
\end{algorithmic}

\section{Complexity Analysis}
\begin{itemize}
    \item \textbf{Time:} $O(n_P \cdot d_P \cdot s_{\max} + n_C \cdot d_C \cdot s_{\max} + s_{\max})
          = O(9 \cdot 4 \cdot 36 + 6 \cdot 6 \cdot 36 + 36) = O(1)$.
    \item \textbf{Space:} $O(s_{\max}) = O(36) = O(1)$.
\end{itemize}

\section{Answer}

\[
\boxed{0.5731441}
\]

\end{document}
